%%
%% This is file `example.tex',jjjjjjjjjjjjjjjjjjjjjjjjjjjj
%% generated with the docstrip utility.
%%
%% The original source files were:
%%
%% coppe.dtx  (with options: `example')
%% 
%% This is a sample monograph which illustrates the use of `coppe' document
%% class and `coppe-unsrt' BibTeX style.
%% 
%% \CheckSum{1416}
%% \CharacterTable
%%  {Upper-case    \A\B\C\D\E\F\G\H\I\J\K\L\M\N\O\P\Q\R\S\T\U\V\W\X\Y\Z
%%   Lower-case    \a\b\c\d\e\f\g\h\i\j\k\l\m\n\o\p\q\r\s\t\u\v\w\x\y\z
%%   Digits        \0\1\2\3\4\5\6\7\8\9
%%   Exclamation   \!     Double quote  \"     Hash (number) \#
%%   Dollar        \$     Percent       \%     Ampersand     \&
%%   Acute accent  \'     Left paren    \(     Right paren   \)
%%   Asterisk      \*     Plus          \+     Comma         \,
%%   Minus         \-     Point         \.     Solidus       \/
%%   Colon         \:     Semicolon     \;     Less than     \<
%%   Equals        \=     Greater than  \>     Question mark \?
%%   Commercial at \@     Left bracket  \[     Backslash     \\
%%   Right bracket \]     Circumflex    \^     Underscore    \_
%%   Grave accent  \`     Left brace    \{     Vertical bar  \|
%%   Right brace   \}     Tilde         \~}
%%
\documentclass[grad,numbers]{coppe}
\usepackage{amsmath,amssymb}
\usepackage[hidelinks]{hyperref}
\usepackage[utf8]{inputenc}
\usepackage[brazil]{babel}
\usepackage[T1]{fontenc}
\usepackage{graphicx}
\usepackage{float}
\usepackage{multirow}
\usepackage{rotating}
\usepackage{pdflscape}
\usepackage{gensymb}
\usepackage{caption}
\usepackage[siunitx, american, plcs]{circuitikz}
\usepackage[siunitx]{circuitikz}
\usepackage{subcaption}
\usepackage{siunitx}
\usepackage{indentfirst}
\usepackage[table,xcdraw]{xcolor}
\makelosymbols
\makeloabbreviations

\begin{document}
  \title{\textbf{Relatório final Letras- Libras}}
  \foreigntitle{Thesis Title}
  \author{\textbf \\  Wesley da Silva Vanderlei (119089795) }{}
  \advisor{Prof.}{João}{Paulo}{D.Sc.}
  
  \date{11}{2025}

  \maketitle
  %\frontmatter

  %\makecatalog

  \tableofcontents

  \printlosymbols

  \printloabbreviations
   \mainmatter
%  \doublespacing

\chapter{Introdução}
 
 \section {Introdução}



\newline
\newline
 






 

\section{Desenvolvimento}
\subsection{Aspectos educacionais}


\newline

\subsection{Aspectos culturais}



\subsection{Importância da história}

\newline
	
\newline

\newline






\section{Conclusões}



 \chapter{Explicação do código}
 \chapter{Manual do Usuário}
 \section{Elementos}
 \subsection{Criação de um elemento}
 Para criar um elemento qualquer(contato,bobina) o usuario deve escrever o seguinte codigo:

 \begin{verbatim}
criar_lemento(nome,tipo,extra)
\end{verbatim}

Onde o nome deve ser o nome dado ao elemento na forma de string, tipo é o tipo do elemento em letras minusculas e extra pode ou nao ser colocado


\chapter{Conexão entre elementos}
\section{Powerail}
\subsection{1 elemento a 1 elemento}


Para fazer a conexão de um elemento a outro elemento é necessario usar o comando fio master(('Nome elemento 1','tipo elemento1'),('Nome elemento 2','Tipo elemento 2'):



 \begin{verbatim}
fio_master(('Contato1','contato'),('Bobina1','bobina'))
\end{verbatim}


\subsection{Varios elementos a 1 elemento}

\begin{verbatim}
fio_master([('Contato1','contato'),...('Contaton','contato')],('Bobina1','bobina'))
\end{verbatim}

  
\end{document}
%% 
%%
%% End of file `example.tex'.
